\section*{章节综合：为一次近似负责}
想估算 $f(1+h)=\mathrm e^{(1+h)^2}$，先问需要多高精度，而不是立即展开。链式法则给出 $f'(x)=2x\mathrm e^{x^2}$，再用乘积法则得到 $f''(x)=(2+4x^2)\mathrm e^{x^2}$。所以一点附近的二阶近似为
\[f(1+h)\approx \mathrm e+2\mathrm e h+3\mathrm e h^2.\]
当 $h=0.01$ 时，预测值为 $1.0203\,\mathrm e$。一次近似只保留 $1.02\,\mathrm e$，两者差 $0.0003\,\mathrm e$，正是二阶修正，而非某种计算器误差。

要声称准确到某位小数，还需给余项界。求得 $f'''(x)=(12x+8x^3)\mathrm e^{x^2}$ 后，在一到 $1.01$ 的区间给它一个上界 $M$，误差即可控制在 $M|h|^3/6$ 内。局部匹配与可证明精度因此是相邻但不同的步骤。

请依次解释内层平方改变输入刻度、外层指数改变输出、二阶导数修正斜率变化、余项约束剩余误差。若能把这四句连起来，导数法则与泰勒公式就不再是分离的计算技巧。
